\documentclass[11pt]{article}
\usepackage[T1]{fontenc}
\usepackage[utf8]{inputenc}
\usepackage{lmodern}
\usepackage[a4paper,margin=1in]{geometry}
\usepackage{microtype}
\usepackage{amsmath,amssymb,amsthm,mathtools}
\usepackage{booktabs,array}
\usepackage[dvipsnames]{xcolor}
\usepackage[colorlinks=true,linkcolor=MidnightBlue,citecolor=MidnightBlue,
            urlcolor=MidnightBlue]{hyperref}
\usepackage{enumitem}
\usepackage{setspace}
\onehalfspacing
\setlist[itemize]{topsep=4pt,itemsep=2pt,leftmargin=1.5em}
\setlist[enumerate]{topsep=4pt,itemsep=2pt,leftmargin=1.8em}

\newtheorem{theorem}{Theorem}[section]
\newtheorem{lemma}[theorem]{Lemma}
\newtheorem{proposition}[theorem]{Proposition}
\newtheorem{corollary}[theorem]{Corollary}
\theoremstyle{definition}
\newtheorem{definition}[theorem]{Definition}
\newtheorem{example}[theorem]{Example}
\theoremstyle{remark}
\newtheorem{remark}[theorem]{Remark}
\newtheorem{observation}[theorem]{Observation}

\newcommand{\DR}{\mathsf{DR}}
\newcommand{\PO}{\mathsf{PO}}
\newcommand{\PP}{\mathsf{PP}}
\newcommand{\PPI}{\mathsf{PPI}}
\newcommand{\EQ}{\mathsf{EQ}}
\newcommand{\RCC}{\{\DR,\PO,\PP,\PPI\}}
\newcommand{\RCCall}{\{\DR,\PO,\EQ,\PP,\PPI\}}
\newcommand{\cl}{\operatorname{cl}}
\newcommand{\Tp}{\operatorname{Tp}}
\newcommand{\tp}{\operatorname{tp}}
\newcommand{\comp}{\operatorname{comp}}
\newcommand{\inv}{\operatorname{inv}}
\newcommand{\TypeSafe}{\operatorname{TypeSafe}}
\newcommand{\ALCIRCC}[1]{\mathcal{ALCI}_{\mathsf{RCC#1}}}

\title{\textbf{A Two-Node Counterexample to Naive Typed Patchwork}\\[0.4em]
\large Witness-menu propagation is load-bearing for $\ALCIRCC{5}$}
\author{Claude (Anthropic, Opus 4.7)\\
\small Prompted by Michael Wessel}
\date{April 2026}

\begin{document}
\maketitle

\begin{center}
\fbox{\parbox{0.92\textwidth}{%
\textbf{Disclaimer.}\;
This document was generated by Claude (Anthropic), prompted by Michael
Wessel.  It corrects a conjecture made in an earlier probe paper by
the same author~\cite{PatchworkProbe}.  The results have not been
independently verified.  We invite scrutiny and corrections.}}
\end{center}

\begin{abstract}
In an earlier probe paper~\cite{PatchworkProbe} we conjectured
(Proposition~7.1) that arc-consistent typed $\mathsf{RCC5}$ networks
with non-empty filtered domains are always realisable by some model of
the $\ALCIRCC{5}$ TBox.  The conjecture, under the weak form of
arc-consistency used in that paper (composition plus type-safety
propagation over existing nodes only), is \emph{false}.  We exhibit a
two-node counterexample: a TBox with seven axioms under which the
network $\{x, z\}$ with natural Hintikka types is arc-consistent but
has no model, because a forced $\exists\PP.D$ witness at $x$ cannot
be related to $z$ by any atomic relation compatible with $z$'s
universals.  The failure is local in the network but non-local in the
TBox: it sits in the gap between ``every edge of $V$ is type-safe''
and ``every pending existential at $V$ is extensibly type-safe.''  We
analyse why the counterexample works, propose the minimal correction
(adding witness-menu propagation to arc-consistency), and observe that
the corrected conjecture essentially reconstructs the split-forest
paper's rank-$d$ witness-menu machinery~\cite{SplitForest}.  The
correction does not affect the decidability claim for $\ALCIRCC{5}$;
it does however sharpen the structural picture and explains why
split-forest's apparent complexity is load-bearing rather than
cosmetic.
\end{abstract}

\section{Introduction}
\label{sec:intro}

The patchwork-augmentation probe paper~\cite{PatchworkProbe}
investigated whether typed arc-consistency is complete for
realisability of $\ALCIRCC{5}$ typed networks.  Three natural
counterexample probes were attempted; each collapsed to local
arc-inconsistency or type-level unsatisfiability, suggesting that
naive typed patchwork might hold.  The paper conjectured this but did
not prove it (Proposition~7.1, stated conjecturally).

This paper provides the counterexample we missed.  It is smaller than
anything we tried: just two nodes, seven TBox axioms, no triangles.
The failure mode we overlooked was not a failure at any arc of the
existing network, but a failure to \emph{extend} the network with a
pending existential witness.  Arc-consistency on $V$, as we defined
it, does not look past $V$; so a TBox that forces a witness
incompatible with some $z \in V$ is not caught.

The correction is conceptually minor---add existential propagation to
arc-consistency---but the consequence is structural: the corrected
propagation is exactly what the split-forest paper's rank-$d$
witness-menus compute.  The simpler patchwork-style argument we were
hoping for in~\cite{PatchworkProbe} does not exist.  Any honest
decidability proof for $\ALCIRCC{5}$ has to do the witness-menu
work, either directly or via an equivalent mechanism.

\paragraph{What this does and does not say.}
This paper corrects a specific proposition from a specific probe
paper.  It does \emph{not}:
\begin{itemize}
  \item Falsify the decidability of $\ALCIRCC{5}$.  The split-forest
    soundness theorem~\cite{SplitForest} and the completeness
    extraction theorem~\cite{CompletenessExtraction} remain
    intact; neither relied on naive arc-consistency.
  \item Falsify any of the three probes in~\cite{PatchworkProbe}.
    Those still reduce to cheap local obstructions under the
    corrected definition of arc-consistency too; the issue is
    elsewhere.
  \item Affect the cover-tree tableau implementation~\cite{CoverTree}.
    CT1--CT5 are a different propagation scheme and do not use naive
    arc-consistency either.
\end{itemize}

What it does say: the probe paper's sketch of ``a simpler alternative
decidability proof via arity-$\leq 3$ CSP completeness''
(\cite{PatchworkProbe}, \S9) is too optimistic.  Witness
extensibility is a genuinely harder constraint than pairwise
type-safety, and it must be handled explicitly.

\section{The weak form of arc-consistency}
\label{sec:weak}

We recall the setup from~\cite{PatchworkProbe}, using identical
notation.  Fix an $\ALCIRCC{5}$ TBox $T$ and concept $C_0$.

\begin{definition}[Weak arc-consistency; \cite{PatchworkProbe}, Def.~4.3]
\label{def:weak-ac}
A \emph{disjunctive typed network} is $(V, D, \tau)$ with
$V$ finite, $D : V \times V \to 2^\RCCall$ satisfying
$\inv$-symmetry, and $\tau : V \to \Tp(C_0, T)$.  Weak
arc-consistency is the fixed point of
\[
  D'(x, z) := D(x, z) \,\cap\,
  \bigcup_{R_1 \in D(x, y)} \comp(R_1, D(y, z))
  \,\cap\, \TypeSafe(\tau(x), \tau(z))
\]
over all $x, y, z \in V$.  A network is \emph{weakly arc-consistent
with non-empty domains} if the fixed point has $D(x, y) \neq
\varnothing$ for every pair.
\end{definition}

This propagation uses only (i) the RCC5 composition table on triples
within $V$ and (ii) pairwise type-safety.  Crucially, it ignores the
existential content $\exists R. D \in \tau(x)$: obligations to
\emph{extend} the network are not propagated back to constraints on
existing nodes.

The probe paper conjectured:

\begin{proposition}[\cite{PatchworkProbe}, Prop.~7.1, conjectured]
\label{prop:false}
For every $\ALCIRCC{5}$ TBox $T$ and satisfiable concept $C_0$,
every weakly arc-consistent typed network with non-empty domains
is realisable by some model of $T \cup \{C_0\}$.
\end{proposition}

We now show Proposition~\ref{prop:false} is false.

\section{The counterexample}
\label{sec:ce}

\begin{example}[Two-node unrealisable network]
\label{ex:two-node}
Let $T$ be the following $\ALCIRCC{5}$ TBox with concept names
$A, B, D$:
\begin{align*}
  \text{(T1)} \quad A &\sqsubseteq \exists \PP. D, \\
  \text{(T2)} \quad A &\sqsubseteq \neg B, \\
  \text{(T3)} \quad B &\sqsubseteq \forall \DR. \neg D, \\
  \text{(T4)} \quad B &\sqsubseteq \forall \PO. \neg D, \\
  \text{(T5)} \quad B &\sqsubseteq \forall \PP. \neg D, \\
  \text{(T6)} \quad B &\sqsubseteq \forall \PPI. \neg D, \\
  \text{(T7)} \quad D &\sqsubseteq \neg B.
\end{align*}
Take $C_0 = A$ and $V = \{x, z\}$ with
\begin{align*}
  \tau(x) &\;\supseteq\; \{A,\, \neg B,\, \neg D,\, \exists \PP. D\}, \\
  \tau(z) &\;\supseteq\; \{B,\, \neg A,\, \neg D,\,
      \forall \DR. \neg D,\, \forall \PO. \neg D,\,
      \forall \PP. \neg D,\, \forall \PPI. \neg D\}.
\end{align*}
$\tau(x)$ and $\tau(z)$ are completed to Hintikka types for $(C_0, T)$
by closure under TBox axioms and Boolean saturation; the completions
are consistent and the choice of $\neg D \in \tau(x)$ is TBox-free
(nothing in $T$ forces $D$ on an $A$-typed node).
\end{example}

\begin{observation}[Arc-consistency]
\label{obs:ac}
The pair $(V, D, \tau)$ is weakly arc-consistent with non-empty
domains, where $D(x, z) = \TypeSafe(\tau(x), \tau(z)) = \RCC$.
\end{observation}

\begin{proof}
$V$ has only two nodes and hence no triangles; composition
propagation is vacuous.  Arc-consistency reduces to the type-safety
filter.  We check $\TypeSafe(\tau(x), \tau(z))$ atom-by-atom.
\begin{itemize}
  \item $\DR$: $\tau(z)$ contains $\forall \DR. \neg D$, requiring
    $\DR$-neighbours of $z$ to have $\neg D$.  $\tau(x)$ has $\neg D$.
    Safe.  $\tau(x)$ has no $\DR$-universal.  Safe.
  \item $\PO$: Analogous.  $\tau(z)$'s $\forall \PO. \neg D$ is
    satisfied by $\tau(x)$'s $\neg D$.  Safe.
  \item $\PP(x, z)$: equivalent to $\PPI(z, x)$.  $\tau(z)$'s
    $\forall \PPI. \neg D$ requires $x$ to have $\neg D$.  Yes.  Safe.
  \item $\PPI(x, z)$: equivalent to $\PP(z, x)$.  $\tau(z)$'s
    $\forall \PP. \neg D$ requires $x$ to have $\neg D$.  Yes.  Safe.
  \item $\EQ$: forces $\tau(x) = \tau(z)$, but $A \in \tau(x)$ and
    $\neg A \in \tau(z)$ (from (T2)).  Not safe.
\end{itemize}
Hence $\TypeSafe(\tau(x), \tau(z)) = \RCC$, which is non-empty.  No
further propagation applies.
\end{proof}

\begin{observation}[Unrealisability]
\label{obs:unrealisable}
There is no model $M \models T$ containing elements $x^M, z^M$ of the
prescribed types.
\end{observation}

\begin{proof}
Suppose $M$ were such a model.  From $\exists \PP. D \in \tau(x)$
there is $y \in M$ with $\PP(x, y)$ and $D \in \tp^M(y)$.  By (T7),
$\neg B \in \tp^M(y)$.  Consider the atomic $\mathsf{RCC5}$ relation
$R := \rho^M(y, z) \in \RCCall$.  We rule out every possibility:
\begin{itemize}
  \item $R = \DR$: $\tau(z)$'s $\forall \DR. \neg D$ forces $\neg D
    \in \tp^M(y)$; but $D \in \tp^M(y)$.  Contradiction.
  \item $R = \PO$: analogous, via $\forall \PO. \neg D$.
  \item $R = \PP$: i.e.\ $y \PP z$ equivalently $z \PPI y$.  By (T5)
    applied at $z$ (which has $\forall \PP. \neg D$): if $\PPI(z,y)$,
    then the $\PP$-neighbour of $z$ is $y$; wait: $\PP(y,z)$ means
    $y$ is a proper part of $z$, equivalently $z$ is the
    proper-containing relation $\PPI$ from $y$'s view, and
    equivalently $y$ is a proper part of $z$.
    From $z$'s perspective, $y$ is a $\PPI$-neighbour
    (i.e.\ $z$'s $\PPI$-successor is $y$).  $\tau(z)$'s
    $\forall \PPI. \neg D$ forces $\neg D \in \tp^M(y)$.  But
    $D \in \tp^M(y)$.  Contradiction.
  \item $R = \PPI$: dually, $z \PP y$, so $y$ is a $\PP$-neighbour of
    $z$; $\tau(z)$'s $\forall \PP. \neg D$ forces $\neg D \in
    \tp^M(y)$.  Contradiction.
  \item $R = \EQ$: $y = z$, forcing $\tp^M(y) = \tp^M(z)$, so
    $\{B, D\} \subseteq \tp^M(y)$.  But (T7) gives
    $D \sqsubseteq \neg B$, contradicting $B \in \tp^M(y)$.
\end{itemize}
No relation is possible.  Hence no model.
\end{proof}

\begin{corollary}
Proposition~\ref{prop:false} is false:
Example~\ref{ex:two-node} is weakly arc-consistent with non-empty
domains but has no model.
\end{corollary}

\section{What went wrong}
\label{sec:analysis}

The counterexample factors through a gap in Definition~\ref{def:weak-ac}.
Weak arc-consistency propagates constraints among \emph{existing}
nodes of $V$ via composition and pairwise type-safety.  It does not
propagate the implicit constraint that every pending
$\exists R. D \in \tau(x)$ must have, for each $z \in V$, some
atomic $R' \in \comp(\inv(R), D(x, z))$ type-safe for a plausible
witness type.

In Example~\ref{ex:two-node}, the pending existential
$\exists \PP. D$ at $x$ requires a witness $y$ with $D \in \tau(y)$.
For such a $y$ and the existing $z$:
\begin{itemize}
  \item $\comp(\PPI, D(x, z)) = \comp(\PPI, \RCC) = \RCCall$ (all
    atomic relations are in the composition).
  \item But $\TypeSafe(\tau(y), \tau(z))$ is empty for every
    $\tau(y)$ with $D$, by the same atom-by-atom argument as in
    Observation~\ref{obs:unrealisable}.
\end{itemize}
The \emph{pending} filter is empty, but weak arc-consistency never
computes the pending filter.  The emptiness is invisible from
within $V$.

\begin{remark}[Where the probe paper's probes differ]
The three probes of~\cite{PatchworkProbe} (forced coincidence,
inverse cycle, yardstick) all involved TBoxes where the obstructions
were \emph{between existing nodes} of the tried network.  Forced
coincidence put the universals between $y$ and $z$ both in $V$;
inverse cycle tied up three or four existing nodes; yardstick forced
type-saturation before any network was constructed.  None of the
probes was designed to probe the existential-witness extensibility
gap specifically.  That gap requires the obstruction to sit on an
edge where one endpoint is inside $V$ and the other is a pending
witness outside $V$.  Example~\ref{ex:two-node} does exactly this.
\end{remark}

\section{The fix: witness-menu-style propagation}
\label{sec:fix}

The minimal correction is to augment Definition~\ref{def:weak-ac}
with an existential-propagation clause.

\begin{definition}[Strong arc-consistency]
\label{def:strong-ac}
$(V, D, \tau)$ is \emph{strongly arc-consistent} if it is weakly
arc-consistent (Def.~\ref{def:weak-ac}) \emph{and} additionally: for
every $x \in V$ and every $\exists R. D \in \tau(x)$, there exists a
Hintikka type $\gamma \in \Tp(C_0, T)$ with $D \in \gamma$ and the
forward-propagation $\{E : \forall R. E \in \tau(x)\} \subseteq
\gamma$, such that for every $z \in V \setminus \{x\}$:
\[
  \comp(\inv(R), D(x, z)) \cap \TypeSafe(\gamma, \tau(z))
  \,\neq\, \varnothing.
\]
\end{definition}

The new clause says: every pending existential has a candidate
witness type whose atomic connection to every existing node is
non-empty.  This is exactly the information a \emph{witness menu} in
the split-forest paper's rank-$d$ quotient
framework~\cite{SplitForest} provides: a target type $\gamma$ and a
set of permitted relations $\comp(\inv(R), D(x, z)) \cap
\TypeSafe(\gamma, \tau(z))$ for each existing node.

\begin{proposition}[Fix; not proven here]
\label{prop:fix}
Strong arc-consistency is complete for realisability of typed
$\ALCIRCC{5}$ networks.
\end{proposition}

We do not prove Proposition~\ref{prop:fix} in this paper.  A proof
would repackage the split-forest soundness argument (König's lemma
on canonical refinements, LCA induction) with strong arc-consistency
replacing the more elaborate rank-$d$ witness-menu propagation.  We
conjecture the two are equivalent modulo bookkeeping, but this is
not trivial; the split-forest paper's rank-$d$ hierarchy handles
several additional concerns (EQ-sync, sibling-interface descriptor
table coherence (C1)--(C5)) that strong arc-consistency in its
simplest form does not.

In particular: Example~\ref{ex:two-node} is caught by strong
arc-consistency.  For $\exists \PP. D$ at $x$: no Hintikka type
$\gamma$ containing $D$ has non-empty $\TypeSafe(\gamma, \tau(z))$.
The existential clause of Definition~\ref{def:strong-ac} fails, and
the network is flagged as unrealisable.

\section{Implications}
\label{sec:impl}

\subsection{For decidability}

Nothing changes.  The split-forest soundness
theorem~\cite{SplitForest} and the completeness extraction
theorem~\cite{CompletenessExtraction} remain intact.  Both use
propagation stronger than weak arc-consistency; neither relied on
Proposition~\ref{prop:false}.  Decidability of $\ALCIRCC{5}$ is
unaffected.

\subsection{For a ``simpler alternative proof''}

The probe paper closed with the sketch of an ambitious theorem:
``every $\ALCIRCC{5}$ constraint on a finite network has arity
$\leq 3$, and arc-consistency is complete for arity-$\leq 3$ CSPs
over $\mathsf{RCC5}$'' (\cite{PatchworkProbe}, \S9).  The
counterexample here refutes the naive version of this claim.

Existentials in ALCI are \emph{non-local} constraints: they force
the existence of additional nodes whose relations to the entire
current network must be simultaneously consistent with the TBox.
Packing this into a CSP of fixed arity requires either (a) an
unbounded-arity encoding, or (b) explicit existential variables with
domain $\Tp(C_0, T)$ and complex binary constraints between them and
each $z \in V$.  Option (b) is essentially what witness menus in the
split-forest paper are.

There is no simpler alternative decidability proof of the form
sketched.  Witness-menu-style existential propagation is the
minimum viable structure.

\subsection{For the undecidability landscape}

The probe paper observed that if typed patchwork held, every known
undecidability reduction would be blocked by a single structural
fact (arity-$\leq 3$ CSP completeness).  Now that the naive typed
patchwork is false, that particular narrative does not go through.
However, the undecidability picture is not reopened: the correct
narrative is merely more nuanced.
\begin{itemize}
  \item Strong arc-consistency (Def.~\ref{def:strong-ac}) still has
    the ``no $k$-ary synthesis for $k \geq 4$'' property as argued
    in~\cite{PatchworkProbe}, \S8.  Existential clauses add binary
    constraints between the pending witness type $\gamma$ and each
    $z \in V$, not genuine $k$-ary constraints.
  \item Strong arc-consistency is what the split-forest soundness
    uses, and the soundness is proven.  So strong arc-consistency
    \emph{is} complete for quotient-derived networks, which suffice
    for decidability.
  \item The structural reason no undecidability reduction has
    succeeded remains: ALCI constructs decompose into node-local,
    arc-local, triple-local, and existential-extension constraints,
    all of which are captured by strong arc-consistency + rank-$d$
    bookkeeping.
\end{itemize}

The shift is: the ``load-bearing'' place in the decidability proof
is not ``$k$-ary CSP completeness'' (which would be false) but
``strong arc-consistency with explicit existential propagation is
complete for quotient-derived typed networks'' (which is what the
split-forest paper actually proves).

\section{Honest assessment}
\label{sec:honest}

This paper is a focused correction to a specific conjecture
(Proposition~\ref{prop:false}).  What we establish:
\begin{enumerate}
  \item A concrete two-node, seven-axiom counterexample
    (Example~\ref{ex:two-node}) that is weakly arc-consistent with
    non-empty domains and has no model.
  \item The counterexample factors through the gap between
    intra-$V$ type-safety (which weak arc-consistency sees) and
    witness-extensibility (which it does not).
  \item The minimal fix (Definition~\ref{def:strong-ac}) adds a
    witness-menu-style propagation clause and is conjectured to be
    complete (Proposition~\ref{prop:fix}), though we do not prove
    it.
  \item The split-forest paper's apparent complexity
    (rank-$d$ witness menus, canonical refinements) is load-bearing:
    witness extensibility must be propagated explicitly.
\end{enumerate}

What we do not establish:
\begin{enumerate}
  \item A proof of Proposition~\ref{prop:fix}.  The substance of that
    proof lies in the split-forest paper; formally deriving strong
    arc-consistency completeness in a simpler presentation would be
    worthwhile but non-trivial.
  \item Any change to the decidability picture for $\ALCIRCC{5}$.
    The probe was aimed at undecidability (via a typed-patchwork
    failure); this counterexample is to a proof sketch, not to the
    decidability itself.
  \item A sharper CSP-theoretic framing.  The arity-$\leq 3$ sketch
    in~\cite{PatchworkProbe}, \S9 would need to be reformulated to
    account for existential extension, and we do not attempt this
    reformulation here.
\end{enumerate}

\paragraph{Methodological note.}  The probe paper's three probes were
carefully structured but all \emph{within a fixed $V$}.  The
counterexample here sits on the boundary between $V$ and its
pending-witness extension---structurally different from anything the
probes tested.  This is the methodological lesson: when designing
adversarial probes against a conjecture, the search space must
include obstructions that cross the boundary between the stated
network and its implicit extensions.  Obstructions confined to
existing nodes are not representative.  This is the same
``distribution blindness'' that Opus~4.7's adversarial review
surfaced in the cover-tree tableau~\cite{Round2Review}: tests
designed within the generator's distribution miss bugs outside it.
A similar principle applies to conjectures.

\begin{thebibliography}{99}

\bibitem{PatchworkProbe}
Claude (Anthropic, Opus 4.7), prompted by M.~Wessel.
\newblock Patchwork preservation under $\ALCIRCC{5}$ augmentation: A
  probe toward undecidability.
\newblock Manuscript, \texttt{papers/patchwork\_augmentation\_ALCIRCC5.pdf},
  April 2026.

\bibitem{SplitForest}
GPT-5.4 (OpenAI), prompted by M.~Wessel.
\newblock Sibling-interface descriptors and split-forest semantics for
  $\ALCIRCC{5}$.
\newblock Manuscript,
  \texttt{papers/trees/sibling\_interface\_descriptors\_ALCIRCC5\_completed\_eqsync\_canonical\_needpatched.pdf},
  April 2026.

\bibitem{CompletenessExtraction}
Claude (Anthropic, Opus 4.6), prompted by M.~Wessel.
\newblock Completeness extraction for $\ALCIRCC{5}$.
\newblock Manuscript,
  \texttt{papers/completeness\_extraction\_ALCIRCC5.pdf}, April 2026.

\bibitem{CoverTree}
Claude (Anthropic, Opus 4.6 and Opus 4.7), prompted by M.~Wessel.
\newblock A cover-tree tableau for $\ALCIRCC{5}$.
\newblock Manuscript,
  \texttt{papers/cover\_tree\_tableau\_ALCIRCC5.pdf}, April 2026.

\bibitem{Round2Review}
Claude (Anthropic, Opus 4.7), commissioned by M.~Wessel.
\newblock Adversarial review of the $\ALCIRCC{5}$ decidability proof
  and cover-tree tableau (Rounds 1--2).
\newblock Manuscripts in \texttt{papers/review*/}, April 2026.

\end{thebibliography}

\end{document}
